% Root proof; independent full literal audits by D and B passed.
\subsection{The eighteen-arrow three-object strong-component cases}

\begin{lemma}[A singleton and a three-state right-zero stalk]
\label{C:right-zero-singleton-three-stalks}
Let $S$ be a finite two-object category, let
$R=\{1,p,q\}$ be the identity-adjoined right-zero two-band,
and let compatible left $R$-sets and right $S$-actions form
the cross sets of a collage with source block $S$ and sink
corner $R$. Suppose the two stalk sizes are one and three.
If all degree-two rational products vanish on every ordinary
finite right $S$-act cone, then a one-dimensional second
rational cohomology group of this collage has zero square.
\end{lemma}
\begin{proof}
Use the natural right-$S$ diagram of left-$R$ action graphs,
as in Lemma~\ref{C:right-zero-action-graph}. The singleton
stalk gives a point. At the three-state stalk, either the graph
is a forest, or it is a point with two loops, or it is a loop
together with an isolated vertex. Indeed a common core of
three elements gives a forest with no edges. A core of two
elements gives one edge, producing either a forest or a loop
and an isolated vertex. A core of one element gives two loops.
The outside edges join the two images under the minimal
retractions.

If that graph is a forest, the actual natural graph-to-component
map is an objectwise homotopy equivalence. The category bar
and its cofiber therefore identify the collage with the
ordinary act cone on the finite component diagram. Its
degree-two products vanish by hypothesis.

In either loop case, double every loop and fix every vertex.
This is a strict natural self-map of the whole $S$-diagram.
An endomorphism at the three-state stalk sends a loop to a
loop or a vertex; maps to or from the singleton stalk fix the
selected vertices and collapse loops. There is no other
nontrivial edge whose parametrization would require a choice.
Thus the degree-two circle maps commute with every actual
diagram map.

Let $q:X\to Y$ be the actual cofiber comparison to the ordinary
component-act cone. The graph-chain kernel has homology only
in degree one, naturally equal to the graph first-homology
module. The self-map is multiplication by two on that module
and the identity on the component diagram. Apply the free
category bar and its homology exact sequence. Exactly as in
Lemma~\ref{C:graph-component-eigenvalues}, on each rational
cohomology group its pullback satisfies
\[
 (f^*-1)(f^*-2)=0,\qquad
 \ker(f^*-1)=\operatorname{im}(q^*).
\]
This argument uses the objectwise natural graph-chain maps,
so it applies to a finite category in place of a monoid.
On the one-dimensional second cohomology the eigenvalue is
one or two. In the first case the class pulls back from $Y$
and has zero square. In the second its square has eigenvalue
four, forbidden by the displayed polynomial. This proves
the assertion.
\end{proof}

\begin{lemma}[Two-state actions of the five-element sphere band]
\label{C:sphere-band-two-state-strong-source}
Let $N$ be the identity-adjoined $2\times2$ rectangular band,
and let $S$ be a finite strongly connected category with
contractible classifying space. Suppose every cross stalk
of an $(N,S)$-collage has two elements, and $H^2(BS;\mathbb Q)=0$.
Then the collage cannot have the rational cohomology ring
of $\mathbb CP^2$.
\end{lemma}
\begin{proof}
If one stalk has an $N$-fixed state, its action category
contains an actual copy of $BN$ over $BN$. The cross-arrow
space in the mixed-nerve pushout therefore surjects onto
$H_2(BN;\mathbb Q)$. Since $H^2(BS;\mathbb Q)=0$, every
degree-two class of the collage is relative to its two
vertex blocks. Their quotient is a suspension, so its
degree-two products vanish.

Otherwise no stalk has an $N$-fixed state. Every minimal
band element acts idempotently on a two-element set, hence
as the identity or as a constant. If $A_{i_0j_0}=1$ for
the images $A_{ij}$ of the minimal actors, the equations
$A_{i_0j_0}A_{kl}=A_{i_0l}$ first make every image depend
only on its column. Then
$A_{ij}A_{i_0j_0}=A_{ij_0}=1$ gives $A_{ij}=1$ for
every pair. This gives fixed states, contrary to the
present case. Thus every minimal actor is constant.
The equations then force its constant value to depend
only on its row. The two rows must give the two distinct
states, or a common fixed state would again result.

Consequently every stalk is canonically the same row
$N$-set, and every map between stalks is the unique
row-preserving bijection. The whole right $S$-action on
this row set is therefore trivial after these canonical
identifications. The cross-arrow space is
$BS\times B(N\ltimes\{0,1\})$. Contractibility of $BS$
and homotopy invariance of the mixed-nerve pushout identify
the collage with the ordinary left $N$-act cone on this
two-element set. Its degree-two products vanish by the
two-state act theorem.
\end{proof}

\begin{theorem}\label{C:eighteen-three-strong-component}
Every finite category with three objects, at most eighteen
arrows, and a nontrivial strongly connected component is
excluded from type $\mathbb CP^2$.
\end{theorem}
\begin{proof}
Assume such a category is a projective plane. It is connected.
A retract object can be deleted, giving a two-object category
with at most fifteen arrows, contrary to the global lower
bound. A wholly strongly connected category is excluded by
the exact strongly connected minimum twenty-three.
Take opposites so that the proper strong pair $S$ is the
source and the remaining corner $N$ the sink. Put
$s=|S|$, $m=|N|$, and let $t$ be the total cross-set size.
Then $s\geq6$, $t\geq2$, and $s+m+t\leq18$.
Projective trivial right modules at $N$ are excluded by
Corollary~\ref{C:strong-pair-projective-endpoint21}.
Thus $m\geq3$ and that module is nonprojective.

If both vertex blocks have zero second rational cohomology,
Lemma~\ref{C:two-block-zero-H2-products} applies.
Suppose first $H^2(BS;\mathbb Q)\ne0$. The categorical Rees
budget gives $s\geq11$. If $t=2$, every stalk is a singleton
and the collage is a join, so its positive products vanish.
For $t\geq3$ the budget leaves $m\leq4$. If both stalks
have size at most two, use
Lemma~\ref{C:small-corner-two-states-per-object}.
The only remaining possibility is
\[
 (s,m,t)=(11,3,4),\qquad
 (|F(a)|,|F(b)|)=(1,3)\text{ or }(3,1).
\]
The nonprojective three-element corner is $R=\{1,p,q\}$
with a right-zero minimal two-band. The strong pair has
trivial minimal group, total row and column counts three,
and only its two identities outside the minimal band.
It therefore has at most two minimal rows at each object.
Theorem~\ref{C:categorical-two-row-act-cups} and the first
lemma exclude this last case.

It remains that $H^2(BS;\mathbb Q)=0$ and
$H^2(BN;\mathbb Q)\ne0$, so $m\geq5$.
A singleton cross stalk supplies an $N$-fixed state and
an actual section over $BN$ in the cross-arrow space.
As in the preceding lemma, every degree-two class is
relative to the two blocks and has zero square. Thus both
stalks have size at least two, $t\geq4$, and $s\leq9$.
The minimal group of the nonretract strong pair is trivial.
If it has one minimal row at each object, the actual
source-block collapse of
Lemma~\ref{C:one-row-source-block-collapse} gives an ordinary
left $N$-act cone with fewer than twenty-two arrows,
contrary to its exact minimum.

The only remaining Rees shape has total row count three,
total column count two, and row counts $(2,1)$ at $(a,b)$.
Its two column counts are both one. The minimal band and
identities cost eight arrows. If every incoming arrow
$a\to b$ is minimal, delete $b$ using
Lemma~\ref{C:one-sided-corner-deletion}; the result has fewer
than seventeen arrows, a contradiction. Otherwise the sole
extra arrow is a nonminimal $a\to b$, and the budget forces
\[
 s=9,\qquad m=5,\qquad |F(a)|=|F(b)|=2.
\]
The five-element corner with positive second rational
cohomology is the identity-adjoined $2\times2$ band.
The strong source has one minimal column per object and
trivial minimal group, so its classifying space is
contractible: the dual one-column collapse applied to $S$
alone supplies the actual equivalence to a point.
The preceding two-state sphere-band lemma now gives the
contradiction. This exhausts every budget case.
\end{proof}
